\documentclass[12pt,]{article}
\usepackage{lmodern}
\usepackage{amssymb,amsmath}
\usepackage{ifxetex,ifluatex}
\usepackage{fixltx2e} % provides \textsubscript
\ifnum 0\ifxetex 1\fi\ifluatex 1\fi=0 % if pdftex
  \usepackage[T1]{fontenc}
  \usepackage[utf8]{inputenc}
\else % if luatex or xelatex
  \ifxetex
    \usepackage{mathspec}
  \else
    \usepackage{fontspec}
  \fi
  \defaultfontfeatures{Ligatures=TeX,Scale=MatchLowercase}
\fi
% use upquote if available, for straight quotes in verbatim environments
\IfFileExists{upquote.sty}{\usepackage{upquote}}{}
% use microtype if available
\IfFileExists{microtype.sty}{%
\usepackage{microtype}
\UseMicrotypeSet[protrusion]{basicmath} % disable protrusion for tt fonts
}{}
\usepackage[margin=1in]{geometry}
\usepackage{hyperref}
\hypersetup{unicode=true,
            pdftitle={Exceedingly Simple Monotone Regression},
            pdfauthor={Jan de Leeuw},
            pdfborder={0 0 0},
            breaklinks=true}
\urlstyle{same}  % don't use monospace font for urls
\usepackage{color}
\usepackage{fancyvrb}
\newcommand{\VerbBar}{|}
\newcommand{\VERB}{\Verb[commandchars=\\\{\}]}
\DefineVerbatimEnvironment{Highlighting}{Verbatim}{commandchars=\\\{\}}
% Add ',fontsize=\small' for more characters per line
\usepackage{framed}
\definecolor{shadecolor}{RGB}{248,248,248}
\newenvironment{Shaded}{\begin{snugshade}}{\end{snugshade}}
\newcommand{\KeywordTok}[1]{\textcolor[rgb]{0.13,0.29,0.53}{\textbf{#1}}}
\newcommand{\DataTypeTok}[1]{\textcolor[rgb]{0.13,0.29,0.53}{#1}}
\newcommand{\DecValTok}[1]{\textcolor[rgb]{0.00,0.00,0.81}{#1}}
\newcommand{\BaseNTok}[1]{\textcolor[rgb]{0.00,0.00,0.81}{#1}}
\newcommand{\FloatTok}[1]{\textcolor[rgb]{0.00,0.00,0.81}{#1}}
\newcommand{\ConstantTok}[1]{\textcolor[rgb]{0.00,0.00,0.00}{#1}}
\newcommand{\CharTok}[1]{\textcolor[rgb]{0.31,0.60,0.02}{#1}}
\newcommand{\SpecialCharTok}[1]{\textcolor[rgb]{0.00,0.00,0.00}{#1}}
\newcommand{\StringTok}[1]{\textcolor[rgb]{0.31,0.60,0.02}{#1}}
\newcommand{\VerbatimStringTok}[1]{\textcolor[rgb]{0.31,0.60,0.02}{#1}}
\newcommand{\SpecialStringTok}[1]{\textcolor[rgb]{0.31,0.60,0.02}{#1}}
\newcommand{\ImportTok}[1]{#1}
\newcommand{\CommentTok}[1]{\textcolor[rgb]{0.56,0.35,0.01}{\textit{#1}}}
\newcommand{\DocumentationTok}[1]{\textcolor[rgb]{0.56,0.35,0.01}{\textbf{\textit{#1}}}}
\newcommand{\AnnotationTok}[1]{\textcolor[rgb]{0.56,0.35,0.01}{\textbf{\textit{#1}}}}
\newcommand{\CommentVarTok}[1]{\textcolor[rgb]{0.56,0.35,0.01}{\textbf{\textit{#1}}}}
\newcommand{\OtherTok}[1]{\textcolor[rgb]{0.56,0.35,0.01}{#1}}
\newcommand{\FunctionTok}[1]{\textcolor[rgb]{0.00,0.00,0.00}{#1}}
\newcommand{\VariableTok}[1]{\textcolor[rgb]{0.00,0.00,0.00}{#1}}
\newcommand{\ControlFlowTok}[1]{\textcolor[rgb]{0.13,0.29,0.53}{\textbf{#1}}}
\newcommand{\OperatorTok}[1]{\textcolor[rgb]{0.81,0.36,0.00}{\textbf{#1}}}
\newcommand{\BuiltInTok}[1]{#1}
\newcommand{\ExtensionTok}[1]{#1}
\newcommand{\PreprocessorTok}[1]{\textcolor[rgb]{0.56,0.35,0.01}{\textit{#1}}}
\newcommand{\AttributeTok}[1]{\textcolor[rgb]{0.77,0.63,0.00}{#1}}
\newcommand{\RegionMarkerTok}[1]{#1}
\newcommand{\InformationTok}[1]{\textcolor[rgb]{0.56,0.35,0.01}{\textbf{\textit{#1}}}}
\newcommand{\WarningTok}[1]{\textcolor[rgb]{0.56,0.35,0.01}{\textbf{\textit{#1}}}}
\newcommand{\AlertTok}[1]{\textcolor[rgb]{0.94,0.16,0.16}{#1}}
\newcommand{\ErrorTok}[1]{\textcolor[rgb]{0.64,0.00,0.00}{\textbf{#1}}}
\newcommand{\NormalTok}[1]{#1}
\usepackage{graphicx,grffile}
\makeatletter
\def\maxwidth{\ifdim\Gin@nat@width>\linewidth\linewidth\else\Gin@nat@width\fi}
\def\maxheight{\ifdim\Gin@nat@height>\textheight\textheight\else\Gin@nat@height\fi}
\makeatother
% Scale images if necessary, so that they will not overflow the page
% margins by default, and it is still possible to overwrite the defaults
% using explicit options in \includegraphics[width, height, ...]{}
\setkeys{Gin}{width=\maxwidth,height=\maxheight,keepaspectratio}
\IfFileExists{parskip.sty}{%
\usepackage{parskip}
}{% else
\setlength{\parindent}{0pt}
\setlength{\parskip}{6pt plus 2pt minus 1pt}
}
\setlength{\emergencystretch}{3em}  % prevent overfull lines
\providecommand{\tightlist}{%
  \setlength{\itemsep}{0pt}\setlength{\parskip}{0pt}}
\setcounter{secnumdepth}{5}
% Redefines (sub)paragraphs to behave more like sections
\ifx\paragraph\undefined\else
\let\oldparagraph\paragraph
\renewcommand{\paragraph}[1]{\oldparagraph{#1}\mbox{}}
\fi
\ifx\subparagraph\undefined\else
\let\oldsubparagraph\subparagraph
\renewcommand{\subparagraph}[1]{\oldsubparagraph{#1}\mbox{}}
\fi

%%% Use protect on footnotes to avoid problems with footnotes in titles
\let\rmarkdownfootnote\footnote%
\def\footnote{\protect\rmarkdownfootnote}

%%% Change title format to be more compact
\usepackage{titling}

% Create subtitle command for use in maketitle
\providecommand{\subtitle}[1]{
  \posttitle{
    \begin{center}\large#1\end{center}
    }
}

\setlength{\droptitle}{-2em}

  \title{Exceedingly Simple Monotone Regression}
    \pretitle{\vspace{\droptitle}\centering\huge}
  \posttitle{\par}
    \author{Jan de Leeuw}
    \preauthor{\centering\large\emph}
  \postauthor{\par}
      \predate{\centering\large\emph}
  \postdate{\par}
    \date{Version 02, march 30, 2017}


\begin{document}
\maketitle
\begin{abstract}
A C implementation of Kruskal's up-and-down-blocks monotone regression
algorithm for use wth .C(), and a comparison with other implementations.
\end{abstract}

{
\setcounter{tocdepth}{3}
\tableofcontents
}
Note: This is a working paper which will be expanded/updated frequently.
All suggestions for improvement are welcome. The directory
\href{http://deleeuwpdx.net/pubfolders/jbkPava}{deleeuwpdx.net/pubfolders/jbkPava}
has a pdf version, the bib file, the complete Rmd file with the code
chunks, and the R source code.

\section{Introduction}\label{introduction}

O no ! Not another monotone regression implementation !! There are
already so many !!!

There is \texttt{isoreg()} in the \texttt{stats} package (R Development
Core Team (2017)), \texttt{gpava()} in \texttt{isotone()} (De Leeuw,
Hornik, and Mair (2009)), and \texttt{pava} in \texttt{Iso} (Turner
(2015)). There is also \texttt{wmonreg()} in \texttt{smacof} (De Leeuw
and Mair (2009), Mair, De Leeuw, and Groenen (2015)) and
\texttt{amalgm()} from De Leeuw (2016). Now \texttt{gpava()} is written
in R, \texttt{amalgm()} calls the Fortran from Cran (1980),
\texttt{pava()} from \texttt{Iso} calls ratfor Fortran, and
\texttt{isoreg()} only does unweighted least squares. I wanted something
in C (because C is not Fortran) which did weighted monotone regression,
and I wanted to use the .C() interface (because I am exceedingly
simple).

Now \texttt{wmonreg} from \texttt{smacof} fits the bill. It was written
in 2014 by Patrick Groenen and Gertjan van den Burg. Same as our
proposed algorithm here, it implements the up-and-down-blocks algorithm
of Kruskal (1964). If we compare computation time then
\texttt{wmonreg()} is the main competitor. But I also wanted code that
was easy to modify for different unimodal loss functions and that
performed relatively uniformly over the range of ``almost in the correct
order'' to ``order completely wrong''.

\section{Algorithm}\label{algorithm}

The best possible description of the algorithm was already given by
Kruskal (1964) (page 127). We copy his recipe, with a slight change of
notation and terminology.

``Our algorithm starts with the finest possible partitions into blocks,
and joins the blocks together step by step until the correct partition
is found. The finest possible partition consists naturally of n blocks,
each containing only a single \(x_i\).

Suppose we have any partition into consecutive blocks. We shall use
\(\overline{x}_b\) to denote the average of the \(x_i\) in block b. If
\(b_- , b, b_+\) are three adjacent blocks in ascending order, then we
call \(b\) up-satisfied if \(\overline{x}_b < \overline{x}_{b_+}\) and
down- satisfied if \(\overline{x}_{b_-} < \overline{x}_b\) . We also
call \(b\) up-satisfied if it is the highest block, and down-satisfied
if it is the lowest block.

At each stage of the algorithm we have a partition into blocks.
Furthermore, one of these blocks is active. The active block may be
up-active or down-active. At the beginning, the lowest block, consisting
of \(x_{min}\), is up-active. The algorithm proceeds as follows. If the
active block is up-active, check to see whether it is up-satisfied. If
it is, the partition remains unchanged but the active block becomes
down-active; if not, the active block is joined with the next higher
block, thus changing the partition, and the new larger block becomes
down-active. On the other hand, if the active block is down-active, do
the same thing but upside-down. In other words, check to see whether the
down-active block is down-satisfied. If it is, the partition remains
unchanged but the active block becomes up-active; if not, the active
block is joined with the next lower block into a new block which becomes
up-active. Eventually this alternation between up-active and down-
active results in an active block which is simultaneously up-satisfied
and down-satisfied. When this happens, no further joinings can occur by
this procedure, and we transfer activity up to the next higher block,
which becomes up-active. The alternation is again performed until a
block results which is simultaneously up-satisfied and down-satisfied.
Activity is then again transferred to the next higher block, and so
forth until the highest block is up-satisfied and down-satisfied. Then
the algorithm is finished and the correct partition has been obtained."

Our implementation \texttt{jbkPava()}, which uses Kruskal's initials to
distinguish it from other pava's, stays as close as possible to this
description. Blocks are C structures, collected in an array. Each block
has a value, a weight, a size, the index of the previous block, and the
index of the the next block.

\begin{Shaded}
\begin{Highlighting}[]
\KeywordTok{struct}\NormalTok{ block \{}
    \DataTypeTok{double}\NormalTok{ value;}
    \DataTypeTok{double}\NormalTok{ weight;}
    \DataTypeTok{int}\NormalTok{ size;}
    \DataTypeTok{int}\NormalTok{ previous;}
    \DataTypeTok{int}\NormalTok{ next;}
\NormalTok{\};}
\end{Highlighting}
\end{Shaded}

At the end the block values are expanded and returned in the original
vector. Space for the blocks is allocated and freed in the routine
(i.e.~it is not under the control of R). Given Kruskal's description of
the up-and-down blocks algorithm, the code in the appendix should be
pretty readable.

\section{Timings}\label{timings}

We apply the six monotone regression methods to vectors of length 10,000
with unit weights, repeating each analysis 100 times. We report the
elapsed time from \texttt{system.time()}.

\subsection{Random}\label{random}

The vectors of length 10,000 are normal random numbers (a different
vector for each of the runs).

\begin{verbatim}
## wmonreg        0.386
\end{verbatim}

\begin{verbatim}
## gpava        280.351
\end{verbatim}

\begin{verbatim}
## isoreg         0.128
\end{verbatim}

\begin{verbatim}
## amalgm         2.550
\end{verbatim}

\begin{verbatim}
## Iso::pava      6.544
\end{verbatim}

\begin{verbatim}
## jbkPava        0.137
\end{verbatim}

Clearly \texttt{jbkPava()} and \texttt{isoreg()} are the clear winners,
although \texttt{isoreg()} has the advantage that it does not have to
take weighted means and keep track of the block weights. Note that
\texttt{gpava()}, written in R, is abysmally slow. The two methods
\texttt{amalgm()} and \texttt{Iso::pava()} that call Fortran routines
are not really competitive. \texttt{wmonreg()} is doing quite well, but
\texttt{jbkPava()} is more than twice as fast.

\subsection{Reversed}\label{reversed}

We apply our six methods, one hundred times each, to the vector
\texttt{10000:1}, which obviously will produce a completely merged
vector with all elements equal to the mean.

\begin{verbatim}
## wmonreg        4.424
\end{verbatim}

\begin{verbatim}
## gpava        425.842
\end{verbatim}

\begin{verbatim}
## isoreg         0.027
\end{verbatim}

\begin{verbatim}
## amalgm         2.373
\end{verbatim}

\begin{verbatim}
## Iso::pava     10.727
\end{verbatim}

\begin{verbatim}
## jbkPava        0.052
\end{verbatim}

In this situtation, in which blocks are never up-satisfied and always
down-satisfied, again \texttt{jbkPava()} and \texttt{isoreg()} are best.
\texttt{wmonreg()} and \texttt{Iso::pava()} lose ground.

\subsection{Ordered}\label{ordered}

We apply our six methods, one hundred times each, to the vector
\texttt{1:10000}, which obviously just returns the vector itself. There
is no merging at all, block are always up-satisfied and down-satisfied,
and we merely loop through the vector.

\begin{verbatim}
## wmonreg        0.026
\end{verbatim}

\begin{verbatim}
## gpava          5.896
\end{verbatim}

\begin{verbatim}
## isoreg        20.162
\end{verbatim}

\begin{verbatim}
## amalgm         0.049
\end{verbatim}

\begin{verbatim}
## Iso::pava      0.011
\end{verbatim}

\begin{verbatim}
## jbkPava        0.042
\end{verbatim}

\texttt{wmonreg()} is best-in-show, which is interesting because in many
applications we expect the vector to be closer to the correct ordering
than to random or reverse ordering. Both Fortran calling routines
perform well. \texttt{isoreg()} falls flat on its face and is even
beaten by \texttt{gpava()}. There is probably some reasonable
explanation for this, but since unweighted \texttt{isoreg()} is not
really in the competition I did not explore this.

\subsection{Conclusion}\label{conclusion}

Both \texttt{pava()} and \texttt{wmonreg()} perform well over the range
of examples. There is an indication that \texttt{pava()} is considerably
better for vectors that are out of order, while \texttt{wmonreg()} may
be better for vectors that are already close to correct. In future
versions of this paper we will try to speed up \texttt{pava()}
performance.

\section{Code}\label{code}

\subsection{R code}\label{r-code}

\begin{Shaded}
\begin{Highlighting}[]
\KeywordTok{dyn.load}\NormalTok{(}\StringTok{"jbkPava.so"}\NormalTok{)}

\NormalTok{jbkPava <-}\StringTok{ }\ControlFlowTok{function}\NormalTok{ (x, }\DataTypeTok{w =} \KeywordTok{rep}\NormalTok{(}\DecValTok{1}\NormalTok{, }\KeywordTok{length}\NormalTok{(x))) \{}
\NormalTok{  h <-}
\StringTok{    }\KeywordTok{.C}\NormalTok{(}\StringTok{"jbkPava"}\NormalTok{,}
       \DataTypeTok{x =} \KeywordTok{as.double}\NormalTok{(x),}
       \DataTypeTok{w =} \KeywordTok{as.double}\NormalTok{ (w),}
       \DataTypeTok{n =} \KeywordTok{as.integer}\NormalTok{(}\KeywordTok{length}\NormalTok{(x)))}
  \KeywordTok{return}\NormalTok{ (h}\OperatorTok{$}\NormalTok{x)}
\NormalTok{\}}
\end{Highlighting}
\end{Shaded}

\subsection{C code}\label{c-code}

\begin{Shaded}
\begin{Highlighting}[]
\PreprocessorTok{#include }\ImportTok{<stdlib.h>}
\PreprocessorTok{#include }\ImportTok{<stdbool.h>}

\KeywordTok{struct}\NormalTok{ block \{}
    \DataTypeTok{double}\NormalTok{ value;}
    \DataTypeTok{double}\NormalTok{ weight;}
    \DataTypeTok{int}\NormalTok{ size;}
    \DataTypeTok{int}\NormalTok{ previous;}
    \DataTypeTok{int}\NormalTok{ next;}
\NormalTok{\};}

\DataTypeTok{void}\NormalTok{ jbkPava (}\DataTypeTok{double}\NormalTok{ *x, }\DataTypeTok{double}\NormalTok{ *w, }\DataTypeTok{const} \DataTypeTok{int}\NormalTok{ *n) \{}
    \KeywordTok{struct}\NormalTok{ block *blocks = calloc ((}\DataTypeTok{size_t}\NormalTok{) * n, }\KeywordTok{sizeof}\NormalTok{(}\KeywordTok{struct}\NormalTok{ block));}
    \ControlFlowTok{for}\NormalTok{ (}\DataTypeTok{int}\NormalTok{ i = }\DecValTok{0}\NormalTok{; i < *n; i++) \{}
\NormalTok{        blocks[i].value = x[i];}
\NormalTok{        blocks[i].weight = w[i];}
\NormalTok{        blocks[i].size = }\DecValTok{1}\NormalTok{;}
\NormalTok{        blocks[i].previous = i - }\DecValTok{1}\NormalTok{; }
\NormalTok{        blocks[i].next = i + }\DecValTok{1}\NormalTok{;   }
\NormalTok{    \}}
    \DataTypeTok{int}\NormalTok{ active = }\DecValTok{0}\NormalTok{;}
    \ControlFlowTok{do}\NormalTok{ \{}
\NormalTok{        bool upsatisfied = false;}
        \DataTypeTok{int}\NormalTok{ next = blocks[active].next; }
        \ControlFlowTok{if}\NormalTok{ (next == *n) upsatisfied = true;}
        \ControlFlowTok{else} \ControlFlowTok{if}\NormalTok{ (blocks[next].value > blocks[active].value) }
\NormalTok{            upsatisfied = true;}
        \ControlFlowTok{if}\NormalTok{ (!upsatisfied) \{}
            \DataTypeTok{double}\NormalTok{ ww = blocks[active].weight + blocks[next].weight;}
            \DataTypeTok{int}\NormalTok{ nextnext = blocks[next].next;}
            \DataTypeTok{double}\NormalTok{ wxactive = blocks[active].weight * blocks[active].value;}
            \DataTypeTok{double}\NormalTok{ wxnext = blocks[next].weight * blocks[next].value;}
\NormalTok{            blocks[active].value = (wxactive + wxnext) / ww;}
\NormalTok{            blocks[active].weight = ww;}
\NormalTok{            blocks[active].size += blocks[next].size;}
\NormalTok{            blocks[active].next = nextnext;}
            \ControlFlowTok{if}\NormalTok{ (nextnext < *n)}
\NormalTok{                blocks[nextnext].previous = active;}
\NormalTok{            blocks[next].size = }\DecValTok{0}\NormalTok{;}
\NormalTok{        \}}
\NormalTok{        bool downsatisfied = false;}
        \DataTypeTok{int}\NormalTok{ previous = blocks[active].previous;}
        \ControlFlowTok{if}\NormalTok{ (previous == -}\DecValTok{1}\NormalTok{) downsatisfied = true;}
        \ControlFlowTok{else} \ControlFlowTok{if}\NormalTok{ (blocks[previous].value < blocks[active].value) }
\NormalTok{            downsatisfied = true;}
        \ControlFlowTok{if}\NormalTok{ (!downsatisfied) \{}
            \DataTypeTok{double}\NormalTok{ ww = blocks[active].weight + blocks[previous].weight;}
            \DataTypeTok{int}\NormalTok{ previousprevious = blocks[previous].previous;}
            \DataTypeTok{double}\NormalTok{ wxactive = blocks[active].weight * blocks[active].value;}
            \DataTypeTok{double}\NormalTok{ wxprevious = blocks[previous].weight * blocks[previous].value;}
\NormalTok{            blocks[active].value = (wxactive + wxprevious) / ww;}
\NormalTok{            blocks[active].weight = ww;}
\NormalTok{            blocks[active].size += blocks[previous].size;}
\NormalTok{            blocks[active].previous = previousprevious;}
            \ControlFlowTok{if}\NormalTok{ (previousprevious > -}\DecValTok{1}\NormalTok{)}
\NormalTok{                blocks[previousprevious].next = active;}
\NormalTok{            blocks[previous].size = }\DecValTok{0}\NormalTok{;}
\NormalTok{        \}}
        \ControlFlowTok{if}\NormalTok{ ((blocks[active].next == *n) && downsatisfied) }\ControlFlowTok{break}\NormalTok{;}
        \ControlFlowTok{if}\NormalTok{ (upsatisfied && downsatisfied) active = next;}
\NormalTok{    \} }\ControlFlowTok{while}\NormalTok{ (true);}
    \DataTypeTok{int}\NormalTok{ k = }\DecValTok{0}\NormalTok{;}
    \ControlFlowTok{for}\NormalTok{ (}\DataTypeTok{int}\NormalTok{ i = }\DecValTok{0}\NormalTok{; i < *n; i++) \{}
        \DataTypeTok{int}\NormalTok{ blksize = blocks[i].size;}
        \ControlFlowTok{if}\NormalTok{ (blksize > }\FloatTok{0.0}\NormalTok{) \{}
            \ControlFlowTok{for}\NormalTok{ (}\DataTypeTok{int}\NormalTok{ j = }\DecValTok{0}\NormalTok{; j < blksize; j++) \{}
\NormalTok{                x[k] = blocks[i].value;}
\NormalTok{                k++;}
\NormalTok{            \}}
\NormalTok{        \}}
\NormalTok{    \}}
\NormalTok{    free (blocks);}
\NormalTok{\}}
\end{Highlighting}
\end{Shaded}

\section*{References}\label{references}
\addcontentsline{toc}{section}{References}

\hypertarget{refs}{}
\hypertarget{ref-cran_80}{}
Cran, G.W. 1980. ``Algorithm AS 149: Amalgamation of Means in the Case
of Simple Ordering.'' \emph{Journal of the Royal Statistical Society,
Series C (Applied Statistics)}, 209--11.

\hypertarget{ref-deleeuw_E_16d}{}
De Leeuw, J. 2016. ``Exceedingly Simple Isotone Regression with Ties.''
\url{http://deleeuwpdx.net/pubfolders/isotone/isotone.pdf}.

\hypertarget{ref-deleeuw_mair_A_09c}{}
De Leeuw, J., and P. Mair. 2009. ``Multidimensional Scaling Using
Majorization: SMACOF in R.'' \emph{Journal of Statistical Software} 31
(3): 1--30.
\url{http://deleeuwpdx.net/janspubs/2009/articles/deleeuw_mair_A_09c.pdf}.

\hypertarget{ref-deleeuw_hornik_mair_A_09}{}
De Leeuw, J., K. Hornik, and P. Mair. 2009. ``Isotone Optimization in R:
Pool-Adjacent-Violators Algorithm (PAVA) and Active Set Methods.''
\emph{Journal of Statistical Software} 32 (5): 1--24.
\url{http://deleeuwpdx.net/janspubs/2009/articles/deleeuw_hornik_mair_A_09.pdf}.

\hypertarget{ref-kruskal_64b}{}
Kruskal, J.B. 1964. ``Nonmetric Multidimensional Scaling: a Numerical
Method.'' \emph{Psychometrika} 29: 115--29.

\hypertarget{ref-mair_deleeuw_groenen_U_15}{}
Mair, P., J. De Leeuw, and P.J.F. Groenen. 2015. ``Multidimensional
Scaling: SMACOF in R.''
\url{http://deleeuwpdx.net/janspubs/2015/notes/mair_deleeuw_groenen_U_15.pdf}.

\hypertarget{ref-R_17}{}
R Development Core Team. 2017. \emph{R: A Language and Environment for
Statistical Computing}. Vienna, Austria: R Foundation for Statistical
Computing.
\href{\%7Bhttp://www.R-project.org\%7D}{\{http://www.R-project.org\}}.

\hypertarget{ref-turner_15}{}
Turner, R. 2015. \emph{Iso: Functions to Perform Isotonic Regression}.
\href{\%7Bhttps://CRAN.R-project.org/package=Iso\%7D}{\{https://CRAN.R-project.org/package=Iso\}}.


\end{document}
